\lecture{26}{Основы теории графов}

\sourcevideo
    {https://www.youtube.com/playlist?list=PLOzRYVm0a65fhKDS8cYIC7YCuVGJ4FOJx}
    {Week 8: Lecture 26: Basics of Graph Theory}

В распределённой оптимизации агент использует собственную функцию и
сообщения, полученные от разрешённых соседей. Структуру обмена задаёт
граф: его вершины соответствуют агентам, а рёбра --- доступным каналам
связи. Ниже вводятся понятия, необходимые для записи консенсусных
ограничений и построения матрицы Лапласа.

\section{Распределённая постановка}

Пусть агент \(i\) знает функцию
\[
    f_i\colon\mathbb R^d\to\mathbb R,
    \qquad
    i=1,\ldots,N.
\]
Централизованная задача имеет вид
\[
    \min_{x\in\mathbb R^d}
    F(x),
    \qquad
    F(x)=\sum_{i=1}^{N}f_i(x).
\]
В машинном обучении локальная функция может быть эмпирическим риском
\[
    f_i(x)
    =
    \frac1{|\mathcal D_i|}
    \sum_{(\xi,y)\in\mathcal D_i}
    \ell(x;\xi,y).
\]
Локальные минимизаторы в общем случае различаются и не совпадают с
минимизатором суммы.

В распределённой формулировке агент хранит копию
\(x_i\in\mathbb R^d\):
\[
\begin{aligned}
    \min_{x_1,\ldots,x_N}\quad
    &\sum_{i=1}^{N}f_i(x_i),
    \\
    \text{при условии}\quad
    &x_1=\cdots=x_N.
\end{aligned}
\]
При выполнении консенсусного ограничения существует общий вектор
\(x\), для которого \(x_i=x\), поэтому задача эквивалентна исходной.
Одна итерация распределённого метода сочетает локальный
оптимизационный шаг, например
\[
    x_i^k-\alpha_k\nabla f_i(x_i^k),
\]
и обмен векторами с соседями. Стоимость алгоритма определяется не
только вычислениями, но и числом коммуникационных раундов, размером
сообщений и задержками.

Передача локальных данных может не требоваться, однако обмен
параметрами и градиентами сам по себе не обеспечивает формальной
конфиденциальности. Для неё нужны дополнительные механизмы, например
защищённая агрегация или дифференциальная приватность.

\section{Графы, рёбра и соседи}

\begin{definition}[Граф]
Графом называется пара
\[
    \mathcal G=(\mathcal V,\mathcal E),
\]
где \(\mathcal V=\{1,\ldots,N\}\) --- множество вершин, а
\(\mathcal E\) --- множество рёбер.
\end{definition}

В неориентированном графе ребро \(\{i,j\}\) разрешает двусторонний
обмен между агентами \(i\) и \(j\). В ориентированном графе ребро
\((i,j)\) разрешает передачу от \(i\) к \(j\), тогда как обратное
ребро может отсутствовать. Во взвешенном графе ребру сопоставляется
вес \(a_{ij}>0\), характеризующий интенсивность взаимодействия,
доверие или коэффициент усреднения. Для неориентированного графа
обычно предполагается
\[
    a_{ij}=a_{ji}.
\]

\begin{definition}[Соседи]
Для неориентированного графа множество соседей вершины \(i\) равно
\[
    \mathcal N_i
    =
    \{j\in\mathcal V:\{i,j\}\in\mathcal E\}.
\]
Для ориентированного графа определяются
\[
    \mathcal N_i^{\mathrm{out}}
    =
    \{j:(i,j)\in\mathcal E\},
    \qquad
    \mathcal N_i^{\mathrm{in}}
    =
    \{j:(j,i)\in\mathcal E\}.
\]
\end{definition}

Агент может непосредственно использовать состояния только своих
соседей. Информация от удалённой вершины поступает через цепочку
промежуточных узлов. Для ориентированных графов необходимо заранее
зафиксировать соглашение о матрице смежности. Далее используется
\[
    a_{ij}>0
    \quad\Longleftrightarrow\quad
    (i,j)\in\mathcal E,
\]
то есть строка \(i\) описывает исходящие рёбра вершины \(i\).

\section{Пути и связность}

\begin{definition}[Путь]
Путём из \(i\) в \(j\) называется последовательность
\[
    i=v_0,v_1,\ldots,v_\ell=j,
\]
в которой каждая соседняя пара соединена ребром. Число \(\ell\)
называется длиной пути. В ориентированном графе направления рёбер
должны совпадать с направлением движения.
\end{definition}

\begin{definition}[Связность]
Неориентированный граф называется связным, если путь существует между
любыми двумя вершинами. Максимальный связный подграф называется
компонентой связности.
\end{definition}

Если граф несвязен, информация не передаётся между его компонентами.
Поэтому агенты могут согласовать значения внутри каждой компоненты,
но в общем случае не могут решить задачу для полной суммы
\(\sum_i f_i\). Связность является минимальным условием глобального
консенсуса в статической неориентированной сети.

Ориентированный граф называется сильно связным, если ориентированный
путь существует для каждой упорядоченной пары вершин. Он называется
слабо связным, если после удаления направлений получается связный
неориентированный граф. Для двустороннего распространения информации
слабой связности обычно недостаточно.

В изменяющейся сети используется последовательность
\[
    \mathcal G^k=(\mathcal V,\mathcal E^k).
\]
Отдельный граф может быть несвязным. Одно из стандартных условий
требует существования \(T\ge1\), для которого объединённый граф
\[
    \left(
        \mathcal V,
        \bigcup_{t=k}^{k+T-1}\mathcal E^t
    \right)
\]
связен для каждого \(k\). Тогда информация распространяется по всей
сети в течение каждого окна длины \(T\).

\section{Консенсус по рёбрам}

Для графа локальное согласование записывается условиями
\[
    x_i=x_j
    \qquad
    \text{для всех }\{i,j\}\in\mathcal E.
\]

\begin{theorem}[Локальные равенства и глобальный консенсус]
Пусть неориентированный граф \(\mathcal G\) связен. Тогда равенства
на всех рёбрах эквивалентны условию
\[
    x_1=\cdots=x_N.
\]
\end{theorem}

\begin{proof}
Глобальный консенсус очевидно влечёт равенство на каждом ребре.
Обратно, возьмём любые вершины \(i\) и \(j\). По связности существует
путь
\[
    i=v_0,v_1,\ldots,v_\ell=j.
\]
На каждом его ребре выполнено
\(x_{v_{r-1}}=x_{v_r}\), поэтому
\[
    x_i=x_{v_1}=\cdots=x_j.
\]
Поскольку пара \(i,j\) произвольна, все локальные переменные равны.
\end{proof}

Следовательно, на связном графе распределённую задачу можно записать
как
\[
\begin{aligned}
    \min_{x_1,\ldots,x_N}\quad
    &\sum_{i=1}^{N}f_i(x_i),
    \\
    \text{при условиях}\quad
    &x_i=x_j,
    \qquad
    \{i,j\}\in\mathcal E.
\end{aligned}
\]
Каждое ограничение связывает только соседних агентов, что делает
возможной полностью децентрализованную реализацию. Если граф несвязен,
такие равенства обеспечивают консенсус лишь внутри компонент.

\section{Матрицы смежности и степеней}

\begin{definition}[Матрица смежности]
Матрицей смежности называется матрица
\[
    A=(a_{ij})_{i,j=1}^{N},
\]
элементы которой кодируют наличие и вес рёбер. Для простого
неориентированного невзвешенного графа
\[
    a_{ij}
    =
    \begin{cases}
        1,& \{i,j\}\in\mathcal E,\\
        0,& \text{иначе}.
    \end{cases}
\]
При отсутствии петель \(a_{ii}=0\), а неориентированность даёт
\[
    A=A^{\mathsf T}.
\]
\end{definition}

\begin{definition}[Степень]
В неориентированном графе степень вершины равна
\[
    d_i=|\mathcal N_i|.
\]
Во взвешенном случае используется
\[
    d_i^{\mathrm w}=\sum_{j=1}^{N}a_{ij}.
\]
Матрица степеней определяется формулой
\[
    D=\operatorname{diag}(d_1,\ldots,d_N).
\]
\end{definition}

Для неориентированного графа степень равна сумме строки и, благодаря
симметрии, сумме столбца:
\[
    d_i=\sum_j a_{ij}=\sum_j a_{ji}.
\]
Для ориентированного графа при принятом соглашении
\[
    d_i^{\mathrm{out}}=\sum_j a_{ij},
    \qquad
    d_i^{\mathrm{in}}=\sum_j a_{ji}.
\]
Соответственно вводятся диагональные матрицы
\(D^{\mathrm{out}}\) и \(D^{\mathrm{in}}\).

Рассмотрим неориентированный граф с рёбрами
\[
\begin{aligned}
    \mathcal E
    =
    \bigl\{
        &\{1,2\},\{2,3\},\{2,4\},
        \{3,4\},\\
        &\{4,5\},\{4,6\}
    \bigr\}.
\end{aligned}
\]
Его степени равны
\[
    (d_1,\ldots,d_6)=(1,3,2,4,1,1),
\]
а матрицы имеют вид
\[
    D
    =
    \begin{pmatrix}
        1&0&0&0&0&0\\
        0&3&0&0&0&0\\
        0&0&2&0&0&0\\
        0&0&0&4&0&0\\
        0&0&0&0&1&0\\
        0&0&0&0&0&1
    \end{pmatrix},
\]
\[
    A
    =
    \begin{pmatrix}
        0&1&0&0&0&0\\
        1&0&1&1&0&0\\
        0&1&0&1&0&0\\
        0&1&1&0&1&1\\
        0&0&0&1&0&0\\
        0&0&0&1&0&0
    \end{pmatrix}.
\]
Суммы строк \(A\) совпадают с диагональными элементами \(D\).

\section{Топология и коммуникационная стоимость}

В полном неориентированном графе
\[
    d_i=N-1,
    \qquad
    |\mathcal E|=\frac{N(N-1)}2.
\]
Он обеспечивает прямой обмен между любыми агентами, но требует
квадратичного числа каналов. Разреженный связный граф уменьшает
стоимость одной итерации, сохраняя возможность распространения
информации по всей сети.

Если агент \(i\) передаёт соседям вектор из \(\mathbb R^d\), объём
его обмена за раунд имеет порядок
\[
    O(d_i d),
\]
а суммарный объём по сети --- порядок
\[
    O(|\mathcal E|d),
\]
с точностью до соглашения о подсчёте двух направлений передачи по
неориентированному ребру. Более плотная сеть может ускорять
распространение информации, но повышает стоимость раунда.

Матрица смежности и матрица степеней образуют матрицу Лапласа
\[
    L=D-A.
\]
Она кодирует структуру рёбер, степени, связность и консенсусное
подпространство. Её спектральные свойства и применение в
распределённых алгоритмах рассматриваются в следующей лекции.
